problem:  
Let \( n\ge 3 \) and \( S^{n+1}\subset\mathbb{R}^{n+2} \) be the round sphere.  
Fix a compact subgroup \( G\subset O(n+2) \) acting isometrically with cohomogeneity \( k\ge 2 \), and let \( Q=S^{n+1}/G \) be the orbit space equipped with its natural Riemannian–orbifold structure.  
Consider the sequence of \(G\)-equivariant min–max values \( \{W_j^G\}_{j\ge1} \) produced by the equivariant Almgren–Pitts construction, i.e. the critical areas of \(G\)-invariant sweepouts of \(S^{n+1}\).  

**Open Problem (Equivariant Weyl Law for the Volume Spectrum):**  
Does there exist a positive constant \( a_G > 0 \) depending only on the group \( G \) (and dimension \(n\)) such that  
\[
\lim_{j\to\infty} j^{-\frac{1}{n+1}} W_j^G = a_G\,?
\]
Equivalently, does the sequence of \(G\)-equivariant min–max widths satisfy a Weyl-type asymptotic law analogous to the one proved by Liokumovich–Marques–Neves in the nonequivariant case, and if so, is the constant \(a_G\) determined by a geometric quantity of the orbit space \(Q=S^{n+1}/G\) (for instance, by the volume or dimensional measure of \(Q\))?

Why is it a "good" problem:  
This question proposes a deep and natural **equivariant generalization of the Weyl law for the volume spectrum**, a cornerstone result in modern geometric analysis.  
The known Weyl law states that in the standard (nonequivariant) setting, the \(j\)-th min–max width of a closed manifold grows asymptotically like \(a(n)\, j^{1/(n+1)}\) for a universal constant \(a(n)\).  
Introducing group symmetry fundamentally alters the variational space—critical cycles are restricted to the \(G\)-invariant class, and the parameter space effectively reduces to the **orbit space \(Q=S^{n+1}/G\)**, which can have singularities and lower dimension.  

This problem therefore asks whether the *asymptotic law of geometric critical values* changes when symmetry collapses configuration space, and if so, how this change reflects the geometry of \(Q\). A positive answer would reveal a new invariant \(a_G\) capturing the “spectral density” of \(G\)-invariant minimal hyperspheres, likely depending on representation-theoretic or curvature data of the orbit structure.  

Resolving it requires uniting:
- **Equivariant Almgren–Pitts theory** (to define and control \(W_j^G\)),  
- **Measure-theoretic analysis of singular orbit spaces**, and  
- **Microlocal asymptotics** adapted to symmetric variational problems.  

The formulation is simple—“do the \(G\)-equivariant min–max widths obey a Weyl law, and what does the constant mean?”—yet an answer would open a new analytic frontier linking symmetry, spectral asymptotics, and geometry.  
It exemplifies the ideal of a *good* problem: elementary in appearance, grounded in established theory, but leading into vast unexplored structure connecting geometric measure theory, analysis on singular quotients, and group representation geometry.